\section*{अध्याय \ref{ch:g6:classification-kinship} --- वर्गीकरण और नातेदारी}

\begin{solution}{exo:g6:classification-kinship:1}
किसी जीव की वह ख़ासियत जो उसके पास होती है और जो जाँची जा सकने लायक़
कही जाती है: \omterm{def:g4:classifying-animals:vertebrate}{रीढ़} है; चार \omterm{def:g1:my-body:parts}{उपांग} हैं; \omterm{def:g1:animals-around-us:covering}{रोम} हैं और वह अपने बच्चों को \omterm{def:g2:animals-grow-up:born}{दूध}
पिलाता है (या: \omterm{def:g1:animals-around-us:covering}{पंख} हैं)।
\end{solution}

\begin{solution}{exo:g6:classification-kinship:2}
हर \omterm{def:g6:classification-kinship:attribute}{अभिलक्षण} के मालिकों का समूह पूरी तरह किसी दूसरे के भीतर बैठता है:
सारे \omterm{def:g1:animals-around-us:covering}{रोम} वालों के चार \omterm{def:g1:my-body:parts}{उपांग} होते हैं, और सारे चार उपांगों वाले \omterm{def:g1:animals-around-us:animal}{जंतुओं}
की \omterm{def:g4:classifying-animals:vertebrate}{रीढ़} होती है --- उलटा कभी नहीं, और कटते हुए कभी नहीं। डिब्बों के
भीतर डिब्बे, और कोई भी दोनों में \omterm{def:g1:my-body:parts}{टाँग} फैलाए नहीं।
\end{solution}

\begin{solution}{exo:g6:classification-kinship:3}
क्योंकि उन्होंने उन्हें किसी साझे पूर्वज से विरासत में पाया: कोई
\omterm{def:g6:classification-kinship:attribute}{अभिलक्षण} किसी कुल-रेखा में नीचे उतरता है, इसलिए उसके मालिक उसी पूर्वज
की संतानें होते हैं --- यानी एक कुल।
\end{solution}

\begin{solution}{exo:g6:classification-kinship:4}
ट्राउट सिर्फ़ \omterm{def:g4:classifying-animals:vertebrate}{रीढ़} वाले बिंदु से आगे खड़ी है। ह्वेल \omterm{def:g4:classifying-animals:vertebrate}{रीढ़}, चार \omterm{def:g1:my-body:parts}{उपांग} और
रोम-और-दूध --- तीनों से आगे खड़ी है।
\end{solution}

\begin{solution}{exo:g6:classification-kinship:5}
ह्वेल और बिल्ली: वे रोम-और-दूध वाला दोराहा साझा करती हैं, जो यहाँ
खींचा गया सबसे नया दोराहा है; ट्राउट सबसे पुराने दोराहे पर ही अलग हो
गई थी।
\end{solution}

\begin{solution}{exo:g6:classification-kinship:6}
वे एक ही जानकारी को दो बार कहते हैं: डिब्बा वही शाखा है, अपनी सारी
टहनियों समेत --- किसी दोराहे से आगे की हर चीज़ उस दोराहे का \omterm{def:g6:classification-kinship:attribute}{अभिलक्षण}
विरासत में पाती है और उसी के डिब्बे में बैठती है।
\end{solution}

\begin{solution}{exo:g6:classification-kinship:7}
रोम-और-दूध एक ही कुल-रेखा में नीचे उतरते हैं --- वे पूर्वजों पर निशान
लगाते हैं। और ``मछली जैसी बनावट'' उसी पानी वाले जीवन-ढंग से अजनबियों पर
भी छापी जा सकती है, जैसा ह्वेल दिखाती है: बनावट उधार ली हुई, कुल कहीं
और। \omterm{def:g4:classifying-animals:classification}{वर्गीकरण} कुलों को छाँटता है, इसलिए वह विरासत पर भरोसा करता है और
उधार लिए लिबास पर शक।
\end{solution}

\begin{solution}{exo:g6:classification-kinship:8}
ट्राउट: सिर्फ़ \omterm{def:g4:classifying-animals:vertebrate}{रीढ़}। मेंढक: \omterm{def:g4:classifying-animals:vertebrate}{रीढ़}, चार \omterm{def:g1:my-body:parts}{उपांग}। बिल्ली: \omterm{def:g4:classifying-animals:vertebrate}{रीढ़}, चार \omterm{def:g1:my-body:parts}{उपांग},
रोम-और-दूध। कलचिड़ी: \omterm{def:g4:classifying-animals:vertebrate}{रीढ़}, चार \omterm{def:g1:my-body:parts}{उपांग}, \omterm{def:g1:animals-around-us:covering}{पंख}। निशानों में ही नत्थी होना
दिख जाता है: हर पंक्ति के निशानों में \omterm{def:g4:classifying-animals:vertebrate}{रीढ़} शामिल है; और आवरण वाला निशान
लगाने वाली हर पंक्ति चार उपांगों पर भी निशान लगाती है।
\end{solution}

\begin{solution}{exo:g6:classification-kinship:9}
वह \omterm{def:g4:classifying-animals:vertebrate}{रीढ़} और चार उपांगों से गुज़रती है, और अपनी सूखी शल्कदार \omterm{def:g1:the-five-senses:sense}{त्वचा} के साथ
अपनी ही शाखा पर निकल जाती है --- यानी \omterm{def:g1:animals-around-us:covering}{रोम} और \omterm{def:g1:animals-around-us:covering}{पंख} वाले दोराहों से आगे
नहीं, उनके बग़ल में।
\end{solution}

\begin{solution}{exo:g6:classification-kinship:10}
\omterm{def:g2:animals-grow-up:born}{दूध} पिलाना \omterm{prop:g6:classification-kinship:kinship}{नातेदारी} बताता है: रोम-और-दूध \omterm{prop:g3:sorting-living-things:groups}{स्तनधारी} कुल की विरासत हैं,
इसलिए चमगादड़ की शाखा बिल्ली के साथ जाती है। उड़ना जीवन-ढंग से उधार लिया
गया है --- डैने \omterm{prop:g3:sorting-living-things:groups}{पक्षी} की शाखा पर और चमगादड़ की शाखा पर अलग-अलग पैदा
हुए, जैसे ह्वेल और ट्राउट में धारा-रेखित बनावट।
\end{solution}

\begin{solution}{exo:g6:classification-kinship:11}
वह उसी का है। प्लैटिपस \omterm{prop:g3:sorting-living-things:groups}{स्तनधारियों} के साझे पूर्वज से उतरा है और कुल की
विरासतें (\omterm{def:g1:animals-around-us:covering}{रोम}, \omterm{def:g2:animals-grow-up:born}{दूध}) लिए रहता है; और उसका \omterm{def:g2:animals-grow-up:born}{अंडा} एक पुरानी ख़ासियत है जिसे
उसकी शुरुआती शाखा ने बचाए रखा जबकि बाक़ी ने उसे खो दिया। कुल पूर्वज और
\emph{सारी} संतानें होता है --- अनोखी बातें भी उसी में; और गट्ठर
(\cref{exo:g4:classifying-animals:11}) उसे पहले ही भीतर वोट दे चुका
है।
\end{solution}

\begin{solution}{exo:g6:classification-kinship:12}
अगर शाखाओं की हर जोड़ी मिलती है, तो \omterm{def:g5:biodiversity:species}{प्रजातियों} को लंबे समय में पीढ़ियों
के साथ बदलना ही होगा --- यानी नए \omterm{def:g6:classification-kinship:attribute}{अभिलक्षण} पैदा होते और विरासत में जाते
रहे होंगे --- ताकि एक पुरखा-रेखा ट्राउट, बिल्ली और गुलबहार में शाखा बना
सके। एक वाक्य में निहितार्थ: सारे \omterm{def:g6:exploring-our-environment:components}{सजीव}, बदलाव के साथ, साझे पूर्वजों से
उतरे हैं। माँगा जाने वाला सबूत: उस बदलाव के निशान --- बीच के रूप, पुरानी
चट्टानों में सुरक्षित अवशेष, और ठीक वहीं मिलती पारिवारिक समानताएँ जहाँ
\omterm{def:g6:classification-kinship:tree}{वृक्ष} भविष्यवाणी करता है (यह आगे के किसी साल का काम है)।
\end{solution}

\begin{solution}{pb:g6:classification-kinship:1}
\textbf{1.} ट्राउट: \omterm{def:g4:classifying-animals:vertebrate}{रीढ़}। मेंढक: \omterm{def:g4:classifying-animals:vertebrate}{रीढ़}, चार \omterm{def:g1:my-body:parts}{उपांग}। छिपकली: \omterm{def:g4:classifying-animals:vertebrate}{रीढ़}, चार
\omterm{def:g1:my-body:parts}{उपांग}, सूखी शल्कदार \omterm{def:g1:the-five-senses:sense}{त्वचा}। कबूतर: \omterm{def:g4:classifying-animals:vertebrate}{रीढ़}, चार \omterm{def:g1:my-body:parts}{उपांग}, \omterm{def:g1:animals-around-us:covering}{पंख}। चूहा: \omterm{def:g4:classifying-animals:vertebrate}{रीढ़}, चार
\omterm{def:g1:my-body:parts}{उपांग}, रोम-और-दूध। डॉल्फ़िन: \omterm{def:g4:classifying-animals:vertebrate}{रीढ़}, चार \omterm{def:g1:my-body:parts}{उपांग}, रोम-और-दूध (\omterm{def:g2:animals-grow-up:born}{दूध} पर चिट की
गवाही; फ़ैसला ख़ासियतें करती हैं, और \omterm{def:g2:animals-grow-up:born}{दूध} ही फ़ैसला करने वाली ख़ासियत
है)।

\textbf{2.} \omterm{def:g4:classifying-animals:vertebrate}{रीढ़} --- छहों \omterm{def:g4:classifying-animals:vertebrate}{कशेरुकी} हैं। सबसे चौड़ी साझेदारी सबसे पुराने
साझे पूर्वज पर निशान लगाती है: उस गहराई पर छहों एक ही कुल हैं।

\textbf{3.} मेंढक, छिपकली, कबूतर, चूहा, डॉल्फ़िन। कबूतर के डैने और
डॉल्फ़िन के मीन-पंख काम की पोशाक में चार \omterm{def:g1:my-body:parts}{उपांग} ही हैं: वही विरासत में
मिले \omterm{def:g1:my-body:parts}{उपांग}, अलग-अलग जीवनों के गढ़े हुए --- \omterm{def:g6:classification-kinship:tree}{वृक्ष} विरासत गिनता है,
पोशाक नहीं।

\textbf{4.} क्योंकि साझा पता साझा कुल नहीं होता: ``पानी में रहता है''
एक \omterm{def:g2:where-animals-live:habitat}{आवास} है, जिसे अजनबी भी उधार ले सकते हैं, न कि शरीर के पूर्वजों का
कोई \omterm{def:g6:classification-kinship:attribute}{अभिलक्षण} --- यानी
\cref{def:g3:sorting-living-things:criterion} का पुराना नियम, अब अपने
कारण के साथ।

\textbf{5.} तना $\rightarrow$ \omterm{def:g4:classifying-animals:vertebrate}{रीढ़} वाला दोराहा (ट्राउट उसके बाद अलग
होती है) $\rightarrow$ चार-उपांग वाला दोराहा $\rightarrow$ आवरण की तीन
शाखाएँ: सूखे \omterm{def:g1:animals-around-us:covering}{शल्क} (छिपकली), \omterm{def:g1:animals-around-us:covering}{पंख} (कबूतर), रोम-और-दूध (चूहा और डॉल्फ़िन
साथ-साथ, सबसे नए दोराहे पर); और मेंढक चार-उपांग वाले दोराहे के बाद, हर
आवरण से पहले अलग हो जाता है।

\textbf{6.} दराज़ में चूहे का सबसे क़रीबी चचेरा डॉल्फ़िन है (वे
रोम-और-दूध वाला दोराहा साझा करते हैं)। और ट्राउट का सबसे क़रीबी सब
बराबर हैं --- वह पहले ही दोराहे पर अलग हो गई थी, इसलिए बाक़ी पाँचों
उसके बराबर दूरी वाले चचेरे हैं।

\textbf{7.} \omterm{prop:g3:sorting-living-things:groups}{उभयचर} --- चार उपांगों वाला, खुली नम \omterm{def:g1:the-five-senses:sense}{त्वचा} वाला, और अपना
कोई आवरण-दोराहा न रखने वाला; उसकी शाखा चार-उपांग वाले दोराहे और आवरणों
के बीच अलग होती है।

\textbf{8.} \omterm{def:g4:classifying-animals:vertebrate}{रीढ़} और चार उपांगों से आगे, सूखी शल्कदार शाखा पर: यानी
\omterm{prop:g4:classifying-animals:classes}{सरीसृप} --- बिना देखे वर्गीकृत, क्योंकि फ़ैसला चिटें नहीं, \omterm{def:g6:classification-kinship:attribute}{अभिलक्षण}
करते हैं।

\textbf{9.} डॉल्फ़िन और चूहा ज़्यादा पास बैठते हैं: रोम-और-दूध एक
विरासती दोराहा है। धारा-रेखित बनावट पानी की वर्दी है, जो असंबंधित
तैराकों को एक जैसी बाँट दी जाती है
(\cref{exo:g4:adapted-to-their-world:11}); उसे गिन लिया जाए तो उधार
लिया लिबास कुल की तरह फ़ाइल हो जाता और नत्थी होना चौपट हो जाता।

\textbf{10.} एक ही साझी पूर्वज-प्रजाति --- यानी कोई शुरुआती चार उपांगों
वाला \omterm{def:g4:classifying-animals:vertebrate}{कशेरुकी} --- जिससे मेंढक, छिपकली, कबूतर, चूहा और डॉल्फ़िन, सब उतरे
हैं, और हर एक अपने चार \omterm{def:g1:my-body:parts}{उपांग} अलग पोशाक में लिए हुए है।

\textbf{11.} यह कि छहों --- और उन्हें खींचती कक्षा भी --- एक ही शाखाओं
वाले कुल के हैं, जो साझे पूर्वजों से बदलाव के साथ उतरा है, और मनुष्य भी
बाक़ियों के बीच रोम-और-दूध वाली शाखा पर हैं।

\textbf{12.} नत्थी होना जीवित दुनिया के बारे में एक तथ्य है ---
\omterm{def:g6:classification-kinship:attribute}{अभिलक्षण} सचमुच डिब्बों के भीतर डिब्बों में बैठते हैं, और यह किसी
फ़ाइल-बाबू का हुक्म नहीं है। \omterm{def:g6:classification-kinship:tree}{वृक्ष} उस तथ्य की व्याख्या करता है: साझे
पूर्वजों से मिली विरासत नत्थी \omterm{def:g6:classification-kinship:attribute}{अभिलक्षण} पैदा किए बिना रह ही \emph{नहीं}
सकती, जबकि महज़ फ़ाइल करना दुनिया की इस मेहरबान सुथरेपन को चमत्कार ही
छोड़ देता।
\end{solution}
